\section{The \wcabench{} Dataset}
\label{sec:dataset}

This section documents the data source, scale, schema, domain-specific
decoding rules, temporal splits and known biases. We follow the datasheet
convention \citep{datasheets}: anything a downstream user needs in order to
judge whether the benchmark is appropriate for their question is stated here
or in the accompanying data card.

\subsection{Source and provenance}
\label{sec:data-source}

All data originate from the official \wca{} public database export
\citep{wca_export}, a versioned, freely downloadable snapshot of every
sanctioned competition. We use the \code{Results Export} format
\textbf{v2.0.2} (snake-case column names, including the \code{result\_attempts}
table) with an export date of 2026-09-21. No private, scraped or synthetic data
are used for headline results; a synthetic generator is provided separately for
offline continuous-integration runs, and any result computed on synthetic data
is labelled as such. The code is released under Apache-2.0; the data remain
the property of the \wca{} and are redistributed only as derived aggregates.

\subsection{Scale}
\label{sec:data-scale}

Table~\ref{tab:scale} summarizes the scale of the built artifacts. The raw
export contains millions of individual timed attempts, which is what makes the
benchmark statistically rich and simultaneously computationally demanding.

These counts are \textbf{this benchmark's snapshot statistics}, not figures
published by \wca{}: they are the exact row counts of the pinned export
snapshot and are reproducible from \code{data/processed/manifest.json} and the
reconciliation report.\footnote{Snapshot metadata: export format version
\textbf{v2.0.2}; export date \textbf{2026-09-21} (readme and
\code{metadata.json}); the official export enumerates 14 tables but publishes
\emph{no} per-table row counts.} The official \wca{} export pages expose only
the format version (v2.0.2), the export date and download sizes (the SQL
archive is $390{,}136{,}209$~bytes and the TSV archive $377{,}258{,}709$~bytes,
as of 2026-09-21). We therefore report both the snapshot counts and the
official qualitative scale, and we caution that the \wca{} database grows
continuously, so a newer export will not match these numbers exactly
(Section~\ref{sec:limitations}).

\begin{table}[t]
  \centering
  \small
  \begin{threeparttable}
  \caption{Scale of the \wcabench{} data artifacts. All row counts are
  \emph{snapshot statistics} of this benchmark's pinned export
  (v2.0.2, export date 2026-09-21).}
  \label{tab:scale}
  \begin{tabular}{lr}
    \toprule
    Quantity & Value \\
    \midrule
    Results (person--event--round records) & 6{,}909{,}454 \\
    Competitors (\code{persons}) & 298{,}551 \\
    Competitions (\code{competitions}) & 18{,}708 \\
    Single attempts (\code{result\_attempts}) & 31{,}847{,}257 \\
    Scramble rows (\code{scrambles}) & 3{,}186{,}380 \\
    Events (active + discontinued) & 22 \\
    \midrule
    Train split (2003--2022) & 3{,}211{,}294 \\
    Validation split (2023--2024) & 1{,}978{,}851 \\
    Test split (2025--2026) & 1{,}719{,}290 \\
    \midrule
    Global attempt-level \DNF{} rate & 3.18\% \\
    Frozen person--event statistics & 612{,}224 \\
    \bottomrule
  \end{tabular}
  \begin{tablenotes}\footnotesize
    \item ``Frozen person--event statistics'' counts the per-(person, event)
          historical summaries computed on the training window and held fixed
          for the whole test window (Section~\ref{sec:protocol}).
  \end{tablenotes}
  \end{threeparttable}
\end{table}

The 21 events comprise the 17 currently sanctioned events (three-by-three
through seven-by-seven, one-handed, blindfolded, fewest-moves, multi-blind,
clock, megaminx, pyraminx, skewb, square-1) plus four historical or
discontinued events (three-by-three feet, magic, master magic and old-style
multi-blind). Event coverage is highly imbalanced: three-by-three dominates by
an order of magnitude, whereas multi-blind, four-blind and five-blind have
comparatively few records. This imbalance is the primary reason for the
stratified reporting in Section~\ref{sec:protocol}.

\subsection{Schema}
\label{sec:data-schema}

The benchmark uses the following relational tables, joined through the keys
shown in Table~\ref{tab:schema}.

\begin{table}[t]
  \centering
  \small
  \begin{threeparttable}
  \caption{Core tables and the joins that define the benchmark's relational
  view.}
  \label{tab:schema}
  \begin{tabular}{@{}lll@{}}
    \toprule
    Table & Key fields & Role in the benchmark \\
    \midrule
    \code{persons} & \code{wca\_id}, name, country, gender & competitor identity \\
    \code{competitions} & \code{id}, date, city, coordinates & temporal anchor \\
    \code{results} & (person, competition, event, round) & labels: best, average, pos \\
    \code{result\_attempts} & result key, attempt number & per-attempt values \\
    \code{scrambles} & round key, group & scramble sequences \\
    \code{events} / \code{formats} / \code{round\_types} & id & scoring semantics \\
    \code{countries} / \code{continents} & iso2 & regional stratification \\
    \code{championships} & competition id & optional competition context \\
    \bottomrule
  \end{tabular}
  \begin{tablenotes}\footnotesize
    \item In v2.0.2 the \code{result\_attempts} table no longer carries
          \code{id}/\code{created\_at}/\code{updated\_at}; pipelines must not
          depend on them.
  \end{tablenotes}
  \end{threeparttable}
\end{table}

\subsection{Domain-specific decoding}
\label{sec:data-decoding}

The decisive feature of \wca{} data is that the meaning of a stored integer
depends on the event's \code{format}. We state the decoding rules explicitly,
because mishandling them corrupts every downstream task.

\paragraph{Timed and count-based values.}
For \code{time} events a positive value is a number of centiseconds
(\code{8653} $= 1{:}26.53$). For the fewest-moves event (\code{number}) a
positive value is a raw move count, and averages are stored as $100\times$ the
mean. Negative and zero sentinels encode status: $-1$ means \DNF{} (did not
finish), $-2$ means \DNS{} (did not start), and $0$ means ``no result''.
\DNF{} attempts participate in \DNF{}-rate statistics but never in an average;
\DNS{} is normally excluded from training; $0$ is treated as missing.

\paragraph{Multi-blind compound encoding.}
The multi-blind event (\code{333mbf}) stores a composite integer whose ten
digits pack the number of solved cubes, the number attempted and the elapsed
time into one value. Two generations exist, distinguished by the leading digit.
Algorithm~\ref{alg:multi} gives the decoder we implement (and its exact
inverse, used for round-trip testing).

\begin{algorithm}[t]
\caption{Decoding of the \wca{} multi-blind compound value.}
\label{alg:multi}
\begin{algorithmic}[1]
\Require positive integer $v$ (10-digit, zero-padded string $s$)
\Ensure solved, attempted, seconds
\State $s \gets \mathrm{zfill}(\mathrm{str}(v), 10)$
\If{$s_0 = \texttt{1}$} \Comment{old format \texttt{1SSAATTTTT}}
    \State $SS \gets \mathrm{int}(s_{1:3})$, $AA \gets \mathrm{int}(s_{3:5})$
    \State solved $\gets 99 - SS$
    \State attempted $\gets AA$
    \State seconds $\gets \mathrm{int}(s_{5:10})$
    \State missed $\gets$ attempted $-$ solved \Comment{requires missed $\ge 0$}
\ElsIf{$s_0 = \texttt{0}$} \Comment{new format \texttt{0DDTTTTTMM}}
    \State $DD \gets \mathrm{int}(s_{1:3})$, $MM \gets \mathrm{int}(s_{8:10})$
    \State seconds $\gets \mathrm{int}(s_{3:8})$
    \State difference $\gets 99 - DD$
    \State solved $\gets$ difference $+$ missed, where missed $\gets MM$
    \State attempted $\gets$ solved $+$ missed
\Else
    \State \textbf{error} unknown version digit
\EndIf
\end{algorithmic}
\end{algorithm}

The stored double-blind score used for ranking is the lower-is-better surrogate
$(100-\text{solved})\times 10^6 + \text{seconds} + \text{missed}$, which
preserves the \wca{} ordering while remaining a single scalar per attempt.

\paragraph{Round-format normalization.}
Several round formats require averaging with outlier trimming. For
\emph{average of 5} (\code{ao5}) an attempt set of five is sorted, the best and
worst attempts are discarded, and the remaining three are averaged; if two or
more attempts are \DNF{}, the whole round is \DNF{}. For \emph{mean of 3}
(\code{mo3}) all three attempts are averaged, and any \DNF{} makes the round
\DNF{}. We reconstruct each round's average from \code{result\_attempts} and
check it against the official \code{results.average}; the consistency rate is
tracked as a data-quality gate. This trimming rule matters because a single
\DNF{} in \code{ao5} is usually absorbed as the discarded ``worst'' attempt,
so the marginal effect of one \DNF{} is small, while a second \DNF{} collapses
the round---a discontinuity that models must represent explicitly rather than
average away.

\paragraph{Scrambles.}
Multi-blind scrambles consist of several newline-separated three-by-three
scrambles; in the TSV export newlines are replaced by the pipe character
(\code{|}). The pipeline splits on \code{|} (or newline), trims, and validates
the resulting sequence length against the event.

\subsection{Temporal splits}
\label{sec:data-splits}

We adopt a strict temporal split rather than a random one. Random splits leak
future information, because adjacent competitions of the same competitor are
strongly correlated; a temporal split instead mirrors the real deployment
setting of predicting the future from the past.

\begin{table}[t]
  \centering
  \small
  \begin{threeparttable}
  \caption{Temporal split of \wcabench{}. The test window is fixed so that the
  main leaderboard remains comparable as new data arrive.}
  \label{tab:splits}
  \begin{tabular}{llr}
    \toprule
    Split & Window & Records \\
    \midrule
    Train & 2003--2022 & 3{,}211{,}294 \\
    Validation & 2023--2024 & 1{,}978{,}851 \\
    Test (frozen) & 2025--2026 & 1{,}719{,}290 \\
    \quad Test-A & 2025 H1 & --- \\
    \quad Test-B & 2025 H2 & --- \\
    \quad Test-C & 2026 H1 & --- \\
    \bottomrule
  \end{tabular}
  \begin{tablenotes}\footnotesize
    \item Per-slice record counts are reported at evaluation time and depend on
          the sampled test set; the sub-slices are used for temporal-robustness
          stratification.
  \end{tablenotes}
  \end{threeparttable}
\end{table}

Each competitor additionally retains a complete, time-ordered competition
history, enabling longitudinal analysis. Because the \wca{} database keeps
growing, the main leaderboard is always computed on the fixed 2025--2026 test
window; later competitions are published as a separate \emph{extended test
set}, following the convention of time-series benchmarks.

\subsection{Known biases}
\label{sec:data-biases}

We list the biases that a responsible user must keep in mind; the same list
appears in the data card.
\begin{itemize}
  \item \textbf{Event imbalance.} Three-by-three records vastly outnumber
        multi-blind, four-blind and five-blind records. Unstratified averages
        are therefore dominated by three-by-three.
  \item \textbf{Regional imbalance.} Access to competitions differs sharply by
        country and continent, so cross-regional generalization claims are
        weak by construction.
  \item \textbf{Class imbalance.} \DNF{} is rare (3.18\% of attempts overall)
        and extremely uneven across events; blindfolded and multi-blind events
        have far higher failure rates.
  \item \textbf{Non-stationarity.} Hardware changes (for example magnetic
        cubes), rule changes and training-method shifts make the time series
        non-stationary; a model fitted on 2003--2022 does not face the same
        distribution in 2025--2026.
  \item \textbf{Self-selection.} Participation in an event is a choice, so
        skill-transfer estimates are confounded; simple correlation
        systematically overstates transfer (Section~\ref{sec:tasks}).
  \item \textbf{Cold start.} Competitors whose first appearance falls inside
        the test window have no usable history; they are flagged and reported
        separately rather than silently merged into the main metric.
\end{itemize}
